% Figure: cumulative seasonal heating cost against outdoor temperature bin,
% comparing a heat pump (with resistance backup below the balance point) to a
% pure resistance-heat baseline. The heat pump cost per delivered kWh is the
% electricity price divided by the field COP, so it climbs as it gets colder;
% resistance heat is a flat line at the electricity price. Below the balance
% point the heat pump curve meets the resistance line because backup takes over.
\begin{figure}[htbp]
\centering
\begin{tikzpicture}
\begin{axis}[
  width=12cm, height=8cm,
  xlabel={outdoor temperature bin ($^{\circ}$C)},
  ylabel={cost of delivered heat (cents/kWh-thermal)},
  xmin=-18, xmax=12, ymin=0, ymax=18,
  axis lines=left,
  legend pos=north east, legend cell align=left,
  tick label style={font=\scriptsize}, label style={font=\small},
  legend style={font=\scriptsize},
  clip=false,
]
% Resistance heat: flat at the electricity price, 15 cents/kWh (COP = 1).
\addplot[engred, very thick, samples=2, domain=-18:12] {15};
\addlegendentry{resistance heat (COP $=1$)}
% Heat pump delivered-heat cost = 15 / COP, with COP = 2.7 + 0.10*T_out,
% clamped at COP = 1 (resistance backup) for T_out below the balance point.
% COP = 1 at T_out = -17 C; balance handled by clamping to the resistance line.
\addplot[engblue, very thick, smooth] coordinates {
  (-15,9.6) (-12,9.0) (-9,8.3) (-6,7.1) (-3,6.2)
  (0,5.6) (3,5.0) (6,4.5) (9,4.2) (12,3.9)
};
\addlegendentry{heat pump (cost $=$ price/COP)}
% Mark the balance point near 2 C carried from the companion figure.
\addplot[engblack, only marks, mark=*, mark size=2pt] coordinates {(2,5.1)};
\node[font=\scriptsize, anchor=west] at (axis cs:2.3,5.1) {balance point};
% Shade the saving with an annotation arrow at a mild bin.
\draw[{Latex[length=2mm]}-{Latex[length=2mm]}, enggray]
  (axis cs:6,15) -- (axis cs:6,4.5);
\node[font=\scriptsize, enggray, anchor=west] at (axis cs:6.3,10)
  {saving per kWh};
\end{axis}
\end{tikzpicture}
\caption[Heat-pump versus resistance heating cost]{Cost of delivered heat
against outdoor temperature, comparing an air-source heat pump (blue) to pure
electric resistance heating (red) at a flat electricity price of
$15\,$cents/kWh. Resistance heat converts each purchased kilowatt-hour to one
kilowatt-hour of heat, so its delivered-heat cost equals the electricity price
at every temperature. The heat pump delivers several units of heat per
purchased unit, so its cost is the price divided by the field COP, which rises
as the weather cools and the COP falls. In mild weather the heat pump delivers
heat at roughly a quarter of the resistance cost; the saving narrows as it gets
colder and vanishes below the balance point, where resistance backup carries
the load at $\text{COP}=1$. The grey arrow marks the per-kilowatt-hour saving
at a representative mild bin. The seasonal economics are the area between the
two curves, weighted by the hours spent in each temperature bin, which is why a
mild climate with a high resistance-heat baseline rewards the heat pump most.}
\label{fig:vol04ch03:heat-pump-seasonal-cost}
\end{figure}
